\phantomsection
\part{級数}

\begin{abstract}
この章では，数列の項を限りなく足し合わせた「級数」を導入し，その収束・発散を部分和の数列の極限（\partref{極限}）として定義する．まず級数の収束に関するコーシーの収束条件を確認し，収束級数の一般項が $0$ に収束することを導く．次に，絶対収束の概念を導入し，絶対収束する級数は収束することを示す．そのうえで，級数の線型性や項のまとめ方に関する性質を確認する．最後に，各項が非負である正項級数を扱い，その収束判定に有効な比較定理を証明する．
\end{abstract}

\draftpart

\phantomsection
\section{級数とその性質}

\phantomsection
\subsection{級数の定義}

\begin{definition}{級数}{級数の定義}
  実数列$(a_n)_{n \in \mathbb{N}}$から，数列$(s_n)_{n \in \mathbb{N}}$を
  \[
    s_0 =a_0 ,\quad s_1 =a_0+a_1,\ldots,s_n=a_0+a_1+\cdots+a_n,\ldots
  \]
  と定義する．この数列$(s_n)_{n \in \mathbb{N}}$を，$a_n$を第$n$項とする無限級数といい，$s_n$をこの級数の第$n$部分和ともいう．もし，
  \[
    \lim_{n \to \infty} s_n =s
  \]
  が存在すれば，この級数は収束するといい，$s$をその和とよび，
  \[
    s=\sum_{n=0}^{\infty} a_n =a_0+a_1+\cdots+a_n+\cdots
  \]
  とかく．
\end{definition}

\begin{theorem}{級数に関するコーシーの収束条件}{級数に関するコーシーの収束条件}
  $\sum_{n=0}^{\infty} a_n$が収束することと，任意の$\varepsilon>0$に対し，ある$n_0 \in \mathbb{N}$が存在して，任意の$m,n \in \mathbb{N}$に対して，
  \[
    m,n \geqq n_0 \limp \abs{a_{m+1}+a_{m+2}+\cdots+a_{n}}<\varepsilon
  \]
  となることは同値である．
\end{theorem}

\begin{leftbar}
  \begin{proof}
    前半の主張は，部分和の数列$(s_n)_{n \in \mathbb{N}}$が収束するということであり，このときコーシーの収束条件により，$(s_n)_{n \in \mathbb{N}}$はコーシー列である．ゆえにただちに主張が従う．
  \end{proof}
\end{leftbar}

\begin{cor}{}{}
  級数$\sum_{n=0}^{\infty} a_n$が収束するとき，$\lim_{n \to \infty} a_n =0$である．
\end{cor}

\begin{leftbar}
  \begin{proof}
    \thref{級数に関するコーシーの収束条件}により，このとき，とくに$\lim_{n \to \infty} \abs{s_{n}-s_{n-1}}=0$であり，$\lim_{n \to \infty} a_n =0$となる．
  \end{proof}
\end{leftbar}

\begin{cor}{}{}
  $\lim_{n \to \infty} a_n \ne 0$のとき，$\sum_{n=0}^{\infty} a_n$は発散する．
\end{cor}

\begin{leftbar}
\begin{proof}
  先の命題の対偶をとればよい．
\end{proof}
\end{leftbar}

\begin{example}{級数の収束・発散の判定}{級数の収束・発散の判定}
  $\lim_{n \to \infty} (-1)^n \ne 0$であるから，$\sum_{n=0}^{\infty} (-1)^n$は発散する．
\end{example}

\phantomsection
\subsection{絶対収束}

\begin{prop}{絶対収束}{絶対収束}
  $\sum_{n=0}^{\infty} \abs{a_n}$が収束すれば，$\sum_{n=0}^{\infty} a_n$も収束する．
\end{prop}

\begin{leftbar}
  \begin{proof}
    \[
      \abs{a_{m+1}+a_{m+2}+\cdots+a_{n}} \leqq \abs{a_{m+1}}+\abs{a_{m+2}}+\cdots+\abs{a_{n}}
    \]
    より従う．
  \end{proof}
\end{leftbar}

\begin{remark}[逆は成り立たない]
  \prref{絶対収束}の逆は成り立たない．
  たとえば，$\sum_{n=1}^{\infty} \frac{(-1)^n}{n}$は収束することが知られているが，各項の絶対値の和$\sum_{n=1}^{\infty} \frac{1}{n}$（調和級数）は発散する．
  すなわち，「収束するが，絶対収束はしない」という級数が存在するのである．
\end{remark}

\begin{definition}{絶対収束}{絶対収束}
  $\sum_{n=0}^{\infty} \abs{a_n}$が収束するとき，$\sum_{n=0}^{\infty} a_n$は絶対収束するという．命題\prref{絶対収束}により，絶対収束する級数は収束する．
\end{definition}

\phantomsection
\subsection{級数の線型性および和の性質}

\begin{prop}{級数の線型性}{級数の線型性}
  $\sum_{n=0}^{\infty} a_n$，$\sum_{n=0}^{\infty} b_n$はともに収束して，それぞれ和は$s$，$t$であるとする．このとき，
  \[
    \sum_{n=0}^{\infty} (a_n \pm b_n) = s \pm t,\quad \sum_{n=0}^{\infty} c a_n =cs
  \]
  が成り立つ．ただし$c$を実数の定数とする．
\end{prop}

\begin{leftbar}
  \begin{proof}
    $\sum_{k=0}^{n-1} a_k = s_n$，$\sum_{k=0}^{n-1} b_k = t_n$とすると，
    \begin{align*}
      \sum_{k=0}^{n-1} (a_k \pm b_k) & = s_n \pm t_n \\
      & \to s \pm t~(n \to \infty)
    \end{align*}
    であるから，前半の主張が証明された．後半については，
    \begin{align*}
      \sum_{k=0}^{n-1} ca_k & = c a_n \\
      & \to cs ~(n \to \infty)
    \end{align*}
    となり，これが証明すべきことであった．
  \end{proof}
\end{leftbar}

\begin{prop}{}{}
  $\sum \abs{a_n}$が収束するとき，$\sum a_n$も収束する．
\end{prop}

\begin{leftbar}
  \begin{proof}
    $\sum_{n=0}^{\infty} \abs{a_n}$が収束するとき，\thref{級数に関するコーシーの収束条件}により，任意の$\varepsilon >0$に対して，ある$n_0 \in \mathbb{N}$が存在して，任意の$m,n \in \mathbb{N}$に対して，
    \[
      m,n \geqq n_0 \limp \abs{a_{m+1}} +\abs{a_{m+2}} + \cdots + \abs{a_n}<\varepsilon
    \]
    である．ここで，絶対値の性質により，
    \[
      \abs{a_{m+1}+a_{m+2}+\cdots + a_n } \leqq  \abs{a_{m+1}} +\abs{a_{m+2}} + \cdots + \abs{a_n}
    \]
    であるから，$\sum_{n=0}^{\infty} a_n$も収束する．これが証明すべきことであった．
  \end{proof}
\end{leftbar}

\begin{prop}{}{}
  $\sum_{n=0}^{\infty} a_n$が収束するとき，その連続する有限個の項の和を項とする級数
  \[
    (a_0+a_1+\cdots+a_{n_0})+(a_{n_0 +1}+a_{n_0 +2}+\cdots+a_{n_1})+(\cdots)+\cdots
  \]
  も収束して，その和は$\sum_{n=0}^{\infty} a_n$の和に等しい．
\end{prop}

\begin{example}{}{}
  級数$ \sum a_n$が収束して和$s$をもつとき，たとえば
  \[
    (a_0+a_1) + (a_2+a_3+a_4)+(a_5+a_6)+ \cdots
  \]
  のように項を括弧でくくって得られる級数，すなわち，
  \[
    b_0 = a_0 + a_1 , \quad b_1 =a_2 + a_3 + a_4 ,\quad b_2 =a_5 + a_6 , \cdots
  \]
  として得られる級数$\sum b_n$に関して考える．

  級数$\sum a_n$，$\sum b_n$の部分和をそれぞれ$s_n$，$t_n$とすると
  \begin{align*}
    & t_0 = a_0 + a_1=s_1  , \\
    & t_1 = (a_0 + a_1) + (a_2 + a_3 + a_4)=s_4 ,\\
    & t_2 = (a_0+a_1)+ (a_2 + a_3+a_4)+(a_5+a_6) = s_6
  \end{align*}
  となり，$(t_n)_{n \in \mathbb{N}}$は$(s_n)_{n \in \mathbb{N}}$の部分列である．

  ゆえに，以下の定理により，
  $\lim_{n \to \infty} s_n=s$ならば$\lim_{n \to \infty} t_n =s$である．
\end{example}

\begin{theorem}{}{}
  数列$(a_n)_{n \in \mathbb{N}}$が$\alpha$に収束すれば，その任意の部分列も$\alpha$に収束する．
\end{theorem}

\begin{example}{}{}
  級数$\sum a_n$の第$n$部分和を$s_n$としたとき，数列$(t_n)_{n \in \mathbb{N}}$を
  \[
    t_n=s_{2n}
  \]
  と定義する．このとき，
  \begin{align*}
    & t_0 = s_0=a_0 , \\
    & t_1 = s_2 = a_0 + a_1 + a_2 \\
    & t_2 = s_4 = a_0 + a_1 + a_2 + a_3 + a_4 , \cdots
  \end{align*}
  となり，$(t_n)_{n \in \mathbb{N}}$は$(s_n)_{n \in \mathbb{N}}$の部分列である．

  ゆえに，先の定理により，
  $\lim_{n \to \infty} s_n=s$ならば$\lim_{n \to \infty} t_n =s$である．
\end{example}

\begin{leftbar}
\begin{proof}
  この新しい級数の部分和の列はもとの$\sum_{n=0}^{\infty} a_n$の部分和の数列$(s_n)_{n \in \mathbb{N}}$の部分列$(s_{n(k)})_{n \in \mathbb{N}}$であるから，この命題が成り立つ．
\end{proof}
\end{leftbar}

\phantomsection
\subsection{正項級数}

\begin{definition}{正項級数}{正項級数}
  任意の$n \in \mathbb{N}$に対して，$a_n \geqq 0$である級数$\sum_{n=0}^{\infty} a_n$を正項級数という．
\end{definition}

\begin{theorem}{}{}
  正項級数$\sum_{n=0}^{\infty} a_n$が収束するためには，部分和の数列$(s_n)_{n \in \mathbb{N}}$が上に有界であることが必要かつ十分な条件である．
\end{theorem}

\begin{leftbar}
  \begin{proof}
    \deref{正項級数}により，$a_n \geqq 0$であるから，部分和$(s_n)_{n \in \mathbb{N}}$について，任意の$n \in \mathbb{N}$に対して，
    \[
      s_n \leqq s_{n+1}
    \]
    となり，$(s_n)_{n \in \mathbb{N}}$は単調増加列である．よって，$(s_n)_{n \in \mathbb{N}}$は\thref{単調収束定理}により$\sup \{s_n \mid n \in \mathbb{N}\}$に収束する．逆に$(s_n)_{n \in \mathbb{N}}$が収束すれば，$(s_n)_{n \in \mathbb{N}}$は有界である．
  \end{proof}
\end{leftbar}

\begin{theorem}{比較定理(1)}{比較定理(1)}
  任意の$n \in \mathbb{N}$に対して，$a_n \geqq 0$，$c_n \geqq 0$，$a_n \leqq c_n$が成り立つとする．
  このとき，$\sum_{n=0}^{\infty} c_n$が収束するならば，$\sum_{n=0}^{\infty} a_n$も収束する．
\end{theorem}

\begin{tleftbar}
\begin{proof}
  $(a_n)_{n \in \mathbb{N}}$，$(c_n)_{n \in \mathbb{N}}$が収束するとき，部分和の数列$(A_n)_{n \in \mathbb{N}}$，$(C_n)_{n \in \mathbb{N}}$は上に有界な単調増加列である．
  よって，
  \[
    a \coloneqq \sup \{ A_n \mid n \in \mathbb{N} \} ,\quad c \coloneqq \sup \{ C_n \mid n \in \mathbb{N} \}
  \]
  が存在し，
  \[
    \lim_{n \to \infty} A_n =a ,\quad \lim_{n \to \infty} C_n =c
  \]
  なので，$\sum_{n=0}^{\infty} a_n$は収束する．これが証明すべきことであった．
\end{proof}
\end{tleftbar}

\begin{theorem}{比較定理(2)}{比較定理(2)}
  任意の$n \in \mathbb{N}$に対して，$a_n \geqq 0$，$d_n \geqq 0$，$d_n \leqq a_n$が成り立つとする．
  このとき，$\sum_{n=0}^{\infty} d_n$が発散するならば，$\sum_{n=0}^{\infty} a_n$も発散する．
\end{theorem}

\begin{tleftbar}
\begin{proof}
  背理法で示す．$\sum a_n$が収束すると仮定する．

  このとき，\thref{比較定理(1)}で示したことから，$\sum d_n$も収束するが，これは$\sum d_n$が発散することに矛盾する．

  よって，先の仮定が誤りであり，$\sum a_n$は発散する．これが証明すべきことであった．
\end{proof}
\end{tleftbar}

\newpage

\begin{theorem}{コーシー積}{コーシー積}
  級数$\sum_{n=0}^{\infty} a_n$と$\sum_{n=0}^{\infty} b_n$がともに絶対収束するとする．
  このとき，級数
  \[
    \sum_{n=0}^{\infty} \sum_{k=0}^{n} a_k b_{n-k}
  \]
  も絶対収束し，
  \[
    \sum_{n=0}^{\infty} \sum_{k=0}^{n} a_k b_{n-k}
    = \left( \sum_{n=0}^{\infty} a_n \right) \left( \sum_{n=0}^{\infty} b_n \right)
  \]
  が成り立つ．
\end{theorem}
